%% ===================================================================
%% Appendix B. Equilibrium Sensitivity and the Phase Transition
%% Proofs of thm:phase_transition, the graph-independent nonnormality
%% bound (obs:eta_bound + rem:eta_relu), and the per-node radius
%% (prop:radius). Labels app:proofs / app:proof_phase / app:proof_eta /
%% app:proof_radius are referenced from the main text and preserved.
%% ===================================================================
\section{Equilibrium Sensitivity and the Phase Transition}
\label{app:proofs}

This section proves the three results that turn the operator of \cref{app:prelim} into a safety
boundary: how the equilibrium moves under a perturbation, how far the worst-case slack between the
certificate and the true threshold can grow, and the per-node radius that the certificate yields. Each
proof restates its claim in brief and then derives it step by step using the notation of
\cref{tab:notation}.

\subsection{Proof of the Phase Transition}
\label{app:proof_phase}

\Cref{thm:phase_transition} states that, under (A1)--(A3), a structural perturbation of size
$\varepsilon=\norm{\delta\Ahat}_F$ passes through three regimes set by the critical budget
$\ecrit=(1-\kappa)/\norm{W}_2$: subcritical ($\varepsilon<\ecrit$, the equilibrium moves boundedly),
critical ($\varepsilon\to\ecrit$, its sensitivity diverges), and supercritical
($\varepsilon\ge\ecrit$, contraction is no longer certified).

\begin{proof}
\emph{Step 1: the subcritical sensitivity.}
On each ReLU linear region $F_\theta$ is affine, and the activation pattern at $\zstar$ is stable for
small $\norm{\delta\Ahat}$ because the set of adjacencies whose equilibrium lands exactly on a region
face has Lebesgue measure zero. The conservative calculus of \citet{bolte2021conservative} therefore
supplies an IFT on each region. Differentiating the fixed-point condition $G(z,\Ahat)=z-F(z,\Ahat)=0$ at
$(\zstar,\Ahat)$ gives the equilibrium response
\begin{equation}
\Delta\zstar=(I-J_z)^{-1}J_A\,\mathrm{vec}(\delta\Ahat)+O(\norm{\delta\Ahat}^2),
\label{eq:ift-expand}
\end{equation}
which is exactly $S\,\mathrm{vec}(\delta\Ahat)$ to first order. The resolvent factor is what makes this
an equilibrium rather than a one-step response: a perturbation injected into the loop returns
attenuated by $J_z$ on each pass, and the total is the geometric sum
$\sum_{j\ge0}J_z^{\,j}=(I-J_z)^{-1}$. Since $\kappa=\norm{J_z}_2<1$ this sum converges with
\begin{equation}
\norm{(I-J_z)^{-1}}_2\le\sum_{j\ge0}\kappa^{\,j}=\frac{1}{1-\kappa},
\label{eq:neumann-bound}
\end{equation}
so taking the leading singular value in \eqref{eq:ift-expand} yields the shift bound
$\norm{\Delta\zstar}_F\le\sigma_1(S)\,\varepsilon+O(\varepsilon^2)$ with
$\sigma_1(S)\le\norm{J_A}_2/(1-\kappa)$, i.e.\ \eqref{eq:shift_bound}.

\emph{Step 2: two contraction certificates, and why only the smaller one is rigorous.}
Perturbing the graph moves the state Jacobian to $J_z'=\diag(\phi')(\Ahat'\otimes W)$, with $\diag(\phi')$
the ReLU mask read at the \emph{perturbed} equilibrium. We ask for which $\varepsilon$ the map $F$ is still
a contraction, so that the equilibrium is unique and the first-order expansion of Step~1 governs.

The safe answer ignores the mask. A ReLU mask is a $0/1$ diagonal projection, so multiplying by it can only
shrink a spectral norm, never enlarge it; hence, for \emph{any} activation pattern,
\begin{equation}
\norm{J_z'}_2=\norm{\diag(\phi')(\Ahat'\otimes W)}_2\ \le\ \norm{\Ahat'}_2\norm{W}_2\ \le\ (\norm{\Ahat}_2+\varepsilon)\norm{W}_2,
\label{eq:Jzp-bound}
\end{equation}
the last step using $\norm{\Ahat'}_2\le\norm{\Ahat}_2+\norm{\delta\Ahat}_2\le\norm{\Ahat}_2+\varepsilon$.
The right-hand side falls below $1$ exactly for $\varepsilon<\eglob=1/\norm{W}_2-\norm{\Ahat}_2$. Because
the argument never inspected the mask, it survives as the perturbation drives the equilibrium across
activation regions; a map of spectral norm below $1$ is a Banach contraction, so for every feasible
$\varepsilon<\eglob$ the perturbed model has a single, \emph{globally unique} equilibrium. This is the
rigorous safe radius $\eglob$.

The trained models are contractive for a gentler reason than \eqref{eq:Jzp-bound} sees. Spectral
normalization caps $\norm{W}_2\le c$, but a normalized adjacency with self-loops has $\norm{\Ahat}_2=1$, so
$\eglob=1/c-1$ is small; what holds these models well inside contraction is the ReLU mask, which keeps the
trained factor $\kappa=\norm{J_z}_2$ far under the worst case $\norm{\Ahat}_2\norm{W}_2$. Reading the
threshold off $\kappa$ gives the larger, data-dependent budget $\ecrit=(1-\kappa)/\norm{W}_2\ge\eglob$:
on a \emph{fixed} activation region the perturbed Jacobian is
$J_z'=J_z+\diag(\phi')(\delta\Ahat\otimes W)$, so $\norm{J_z'}_2\le\kappa+\varepsilon\norm{W}_2$, which
reaches $1$ at $\varepsilon=\ecrit$.

The phrase ``fixed activation region'' is where rigor stops. Over a finite budget the equilibrium does not
stay on one region; it crosses finitely many ReLU faces before $\ecrit$, and each crossing can switch on a
dormant unit and add a term the single mask $\diag(\phi')$ omits, so $\kappa+\varepsilon\norm{W}_2$ is not a
valid bound on $\norm{J_z'}_2$ across the whole ball; a free choice of mask can push $\norm{J_z'}_2$ past
$1$ well before $\ecrit$. The two budgets coincide, $\ecrit=\eglob$, only in the all-active case (A4), where
$\kappa=\norm{\Ahat}_2\norm{W}_2$ and no region is crossed. For $\varepsilon<\ecrit$ we therefore claim only
what the implicit function theorem gives locally: the equilibrium continues as a unique smooth branch
wherever $I-J_z$ is invertible, applied region by region through the conservative calculus of
\citet{bolte2021conservative}, together with the first-order bound of Step~1. Whether that branch stays the
unique \emph{global} equilibrium all the way to $\ecrit$ is the empirical regularity of \cref{rem:obs_o1}
(the measured break lies above $\ecrit$), whose two scaling gaps remain open; it holds across our suite,
but we record it as observation, not theorem.

\emph{Step 3: the critical divergence and why $\ecrit$ is one-sided.}
For any matrix $M$ the resolvent $(I-M)^{-1}$ has eigenvalues $(1-\lambda_i(M))^{-1}$, so
\begin{equation}
\norm{(I-M)^{-1}}_2\ \ge\ \rho\!\big((I-M)^{-1}\big)=\frac{1}{\min_i\abs{1-\lambda_i(M)}},
\label{eq:resolvent-lb}
\end{equation}
which blows up when an eigenvalue of $M$ approaches $1$. The rate depends on the geometry of $J_z'$.
When $J_z'$ is normal with a dominant real-positive (Perron) eigenvalue equal to $\norm{J_z'}_2$, the
case for symmetric $\Ahat$ and $W$, the spectral gap to $1$ is
$\min_i\abs{1-\lambda_i}=1-\norm{J_z'}_2=\norm{W}_2(\ecrit-\varepsilon)$ using the within-region
linearization $\norm{J_z'}_2=\kappa+\varepsilon\norm{W}_2$, giving the $\Omega\!\big(1/(\ecrit-\varepsilon)\big)$ divergence. The certificate is only
one-sided because $\norm{J_z'}_2\to1$ need not drive the gap to zero: $M=\diag(-s,0)$ has
$\norm{M}_2=s$ yet $\norm{(I-M)^{-1}}_2=1$, a dominant eigenvalue at $-s$ that stays far from $1$. So
$\ecrit$ lower-bounds the true divergence threshold, and the slack is governed by how far $J_z'$ is from
normal, the nonnormality $\eta$ of \cref{obs:eta_bound}.

\emph{Step 4: supercritical.}
For $\varepsilon\ge\ecrit$ the bound \eqref{eq:Jzp-bound} no longer forces $\norm{J_z'}_2<1$, the Banach
contraction certificate fails, and the iteration may converge to a different fixed point, oscillate, or
diverge.
\end{proof}

\subsection{A graph-independent bound on the nonnormality}
\label{app:proof_eta}

The slack $\eta$ in Step 3 could in principle grow with the graph; the next observation shows it does
not, in the all-active case it is set by $W$ alone.

\begin{observation}[Graph-independent nonnormality bound, all-active case]
\label{obs:eta_bound}
For $J_z=\diag(\phi')(\Ahat\otimes W)$ with $\Ahat$ symmetric and all units active ($\phi'\equiv1$), the
nonnormality obeys $\eta\le\kappa(V_W)$, where $W=V_W D_W V_W^{-1}$ and
$\kappa(V_W)=\norm{V_W}_2\norm{V_W^{-1}}_2$ is the eigenvector conditioning of $W$, independent of
$\Ahat$.
\end{observation}

\begin{proof}
Under (A4), $J_z=\Ahat\otimes W$, whose eigenvalues are the products $\lambda_i(\Ahat)\lambda_j(W)$ and
whose eigenvectors are the Kronecker products of those of $\Ahat$ and $W$ \citep{stewart1990matrix}.
Symmetric $\Ahat$ diagonalizes in an orthonormal basis, so its eigenvector matrix has conditioning
$\kappa(U_{\Ahat})=1$ and contributes nothing to the joint conditioning; the joint eigenvector matrix
therefore has condition number $\kappa(V_W)$. For any diagonalizable $M=VDV^{-1}$ the resolvent obeys
$\norm{(I-M)^{-1}}_2\le\kappa(V)/(1-\rho(M))$, and the slack relative to the spectral-radius rate
$1/(1-\rho(M))$ is exactly $\kappa(V)$. Applied to $J_z$ this gives $\eta\le\kappa(V_W)$.
\end{proof}

\begin{remark}[Empirical extension to general ReLU patterns]
\label{rem:eta_relu}
The bound assumes $\phi'\equiv1$. For general ReLU patterns the diagonalization argument no longer
applies exactly, but the measured nonnormality stays moderate, $\eta\in[1.19,2.47]$ across the suite
(\cref{app:ablations}), tracking $\kappa(V_W)$ closely. Graph structure does not amplify it; the modest
$\eta$ traces to the spectral-norm regularization of $W$, exactly the quantity (A2) controls.
\end{remark}

\subsection{Proof of the Per-Node Radius}
\label{app:proof_radius}

\Cref{prop:radius} gives, for a node $v$ classified as $y_v$ with positive margins
$m_v^{(c)}=f_{y_v}-f_c$ to every competitor $c$, the certified radius
$r_v=\min_{c\ne y_v}m_v^{(c)}/\norm{(W_{y_v}-W_c)S_v}_2$, below which the prediction cannot change.

\begin{proof}
A perturbation of size $\varepsilon$ moves the node's equilibrium by
$\Delta\zstar_v\approx S_v\,\mathrm{vec}(\delta\Ahat)$ (Step 1 of \cref{app:proof_phase}). Because the
head is linear, each margin moves by $(W_{y_v}-W_c)\Delta\zstar_v$, and Cauchy--Schwarz bounds this by
the operator norm of the composed map times the perturbation size,
\begin{multline}
\abs{\Delta m_v^{(c)}}=\abs{(W_{y_v}-W_c)\,S_v\,\mathrm{vec}(\delta\Ahat)}\\
\le\norm{(W_{y_v}-W_c)S_v}_2\,\norm{\delta\Ahat}_F.
\label{eq:margin-move}
\end{multline}
The prediction of $v$ survives as long as every margin stays positive. The margin to competitor $c$
first reaches zero at $\norm{\delta\Ahat}_F=m_v^{(c)}/\norm{(W_{y_v}-W_c)S_v}_2$, and taking the binding
(smallest) value over competitors gives
\begin{equation}
r_v=\min_{c\ne y_v}\frac{m_v^{(c)}}{\norm{(W_{y_v}-W_c)S_v}_2}.
\label{eq:radius-restate}
\end{equation}
The form is a margin divided by how fast the worst perturbation spends it: the denominator is the
product of two amplifiers, the decision-boundary gap $W_{y_v}-W_c$ and the equilibrium sensitivity
$S_v$, so a node earns a large radius exactly when it has a wide margin and a low structural
sensitivity. Minimizing over all competitors measures the distance to the nearest first-order decision
boundary, so the guarantee holds even if the runner-up class changes along the way; substituting the
constrained $S_{c,v}$ for $S_v$ tightens the denominator and yields the larger constrained radius.
\end{proof}
